Analisando limites:
\begin{itemize}
\item $p\to0$: então $w/R\to0$, o que faz sentido porque sem pressão não há deflexão.
\item $E\to\infty$: então $w/R\to0$, pois material muito rígido não deforma.
\item $R\to0$: a relação continua adimensional e pequena se $h$ também for pequena; o comportamento permanece consistente.
\item $h\to0$: então $w/R$ pode crescer, indicando que paredes muito finas se deformam mais facilmente.
\end{itemize}
Esse comportamento qualitativo está de acordo com a intuição física. 